\chapter{Atomic Energies, Li through Rn}\label{ch:atoms}
\index{atomic energies, Li--Rn}\index{shell structure}\index{sub-shell, (4+4)}\index{Aufbau ordering}\index{Neumann boundary condition}\index{periodic table}

% Material to be added later from prel book (realquantum2.tex):
%   - chapters/sphericalsym.tex (computations in spherical symmetry — original
%     1D radial-equation treatment of atoms in the 2017 draft)

\begin{quote}
\itshape Would it not be pre-established harmony of a peculiar sort if the classical-epoch researchers, those who, as we hear today, had no idea of what measuring truly is, had unwittingly gone on to give us as legacy a guidance scheme revealing just what is fundamentally measurable for instance about a hydrogen atom?\\
\normalfont\hfill --- E.~Schr\"odinger, \emph{Die gegenw\"artige Situation in der Quantenmechanik}, 1935.
\end{quote}
\bigskip

Part~III opens with the parameter-free Level-1 atomic model and its validation against the only thing it has to predict: observed total atomic energies. If Level 1 does not reproduce these to acceptable accuracy, the rest of the hierarchy --- Level 2, Level 3, Level 4 --- has nothing to stand on. If it does, every higher-level architecture is anchored to a fixed reference point.

This chapter reports the Level-1 results across the periodic table from Lithium ($Z = 3$) to Radon ($Z = 86$). The agreement with observation is approximately 1\% across the range, with the largest deviations at heavy atoms, where the total energy is a large number assembled from many shells and the simplified computation --- each sub-shell homogenised to spherical symmetry, the atom solved in reduced spatial detail --- accumulates the most error. The agreement is achieved with no fitted parameters: the only input is the nuclear charge.

\section{The parameter-free Level-1 atomic model}\label{sec:level1-model}

The Level-1 model of Section~\ref{sec:level1} is, at full resolution: a point kernel of charge $+Z$ at the origin surrounded by $N$ one-electron wave functions $\psi_i$ on non-overlapping shell-shaped domains $\Omega_i$, with the configuration determined by minimum-energy packing under the Bernoulli free-boundary condition.

The only inputs are:
\begin{itemize}
\item The nuclear charge $Z$.
\item The number of electrons $N$, taken equal to $Z$ for a neutral atom.
\item Computational mesh size, typically $h \sim 0.1$ a.u.\ for minimal but functional resolution.
\end{itemize}
There is no orbital basis, no exchange--correlation functional, no fitted constants, no choice of orbital angular momentum to populate. The variational principle picks the ground-state configuration on its own.

In particular, the model does not begin by writing down the standard hydrogenic orbitals $1s$, $2s$, $2p$ and filling them in Aufbau order. It begins by adding electrons one at a time, each into the lowest-energy spatial region available given the non-overlap constraint and the kernel attraction, and lets the variational principle decide what the ground state looks like.

\section{Emergent shell structure from packing}\label{sec:packing}

The Level-1 model builds up the Periodic Table by sequentially adding electrons to a nucleus of charge $Z$, each electron occupying its own non-overlapping spatial domain $\Omega_i$ with continuity of $\Psi$ across joint boundaries and homogeneous Neumann at the free boundaries. Each electron takes on a width determined by proximity to the kernel while competing for space with the others --- a \emph{balloon-packing} problem in which the kinetic-energy term $K(\Psi)$ from~\eqref{eq:K} acts as a compression cost.

The resulting shell structure is therefore an emergent property of minimum-energy packing, not a presupposed orbital ansatz. This contrasts with StdQM, where shell occupancies follow from Aufbau filling of the $s$, $p$, $d$ orbitals derived from the hydrogen atom.

\paragraph{Shell sequence.}
RealQM packs electrons into an expanding sequence of non-overlapping spherical shells $S_n$ centred at the kernel, $n = 1, 2, 3, \ldots$. Each $S_n$ is split azimuthally into two half-shells, and each half-shell is further partitioned angularly into $n \times n$ electron domains, giving
\[
\text{capacity of shell } S_n \;=\; 2n^2 \text{ electrons.}
\]
This yields the familiar sequence
\[
2,\ 8,\ 18,\ 32,\ \ldots
\]
which is repeated in the Periodic Table periods $2, 8, 8, 18, 18, 32, 32$.

\paragraph{Two-valuedness of the innermost shell.}
For $Z > 1$ the innermost shell $S_1$ already contains \emph{two} half-shell electron densities meeting at a common plane through the kernel, a structure inherited from the Helium atom. This emergent splitting of $S_1$ into two non-overlapping half-spaces accounts for the ``two-valuedness'' which in StdQM has to be supplied externally as the Pauli exclusion principle: in RealQM the non-overlap constraint between distinct electron territories does the same work, and the Pauli principle is not separately needed.

\section{Sub-shell structure}\label{sec:subshells}

RealQM further predicts that filled 8-shells are radially divided into two 4-shells, giving filled-shell decompositions of the form
\[
(2) \;+\; (4{+}4) \;+\; (4{+}4) \;+\; \ldots
\]
This is the central architectural prediction of Level 1 beyond the basic shell counts. The standard $s/p$ distinction of an L-shell (one $s$ orbital, three $p$ orbitals) is replaced by a radial split into two sub-shells of four electrons each. Both pictures yield the same shell occupancy of 8; they differ in how that 8 is internally organised.

The radial sub-shell split is what gives Level 1 quantitative agreement with observation. With Neon at the proper $(2) + (4{+}4)$ structure the computed energy is $-132$~Ha against the observed $-128.5$~Ha. With Neon forced into a single 8-electron outer shell, $(2) + (8)$ without the $(4{+}4)$ subdivision, the computed energy drops to $-153$~Ha, far below observation. The diagnostic is decisive: \emph{Level 1 with sub-shells matches experiment to about 3\%; without sub-shells the energy is wrong by about 19\%}.

\paragraph{Why sub-shells: the Neumann widening argument.}
The mechanism is the homogeneous Neumann condition acting between adjacent radial sub-shells. The Bernoulli condition $\partial\psi_i / \partial n = 0$ at the free boundary widens each shell relative to a free atom: with a flat normal derivative, the wave function cannot decay rapidly across the boundary, and the shell extends radially further than it would under a Dirichlet condition. This widening raises the kinetic-energy cost $K(\Psi)$.

The variational principle therefore prefers configurations that split this cost across two narrower sub-shells rather than concentrate it in a single wider shell. A radially split $(4{+}4)$ configuration distributes the Neumann widening across two boundaries (between the inner 4 and outer 4 sub-shells, and between the outer 4-sub-shell and the next shell) and keeps each individual sub-shell narrow. The single $(8)$ alternative has to accommodate all 8 electrons in one shell whose outer boundary is further from the kernel and whose inner boundary is the same as the $S_1$ outer boundary; the resulting wider shell pays a larger kinetic cost.

Sub-shell structure is thus an \textbf{emergent prediction of the variational principle, not a free choice}. The same Bernoulli condition that closes the multiphase equation predicts a finer radial structure than the standard $s/p$ description, and the predicted structure gives the right energies.

\section{Validation: Li through Rn}\label{sec:li-rn}

Computed ground-state energies for atoms across the periodic table, with each sub-shell homogenised to spherical symmetry in the simplified Atom Simulator code, reproduce observed atomic total energies to within $\sim$1\% for elements as heavy as Radon. Representative results in Hartree are shown in Table~\ref{tab:atoms}.

\begin{table}[h]
\centering
\small
\renewcommand{\arraystretch}{1.10}
\caption{Level-1 RealQM total energies vs observation for representative atoms across the periodic table, with the corresponding shell--sub-shell decomposition.}\label{tab:atoms}
\begin{tabular}{l l r r r}
\toprule
\textbf{Atom} & \textbf{Shell structure} & \textbf{RealQM $E$ (Ha)} & \textbf{Observed (Ha)} & \textbf{Deviation} \\
\midrule
Carbon   & $(2) + (2{+}2)$                                  & $-38.2$    & $-37.7$    & $+1.3\%$ \\
Neon     & $(2) + (4{+}4)$                                  & $-132.4$   & $-128.5$   & $+3.0\%$ \\
Argon    & $(2) + (4{+}4) + (4{+}4)$                        & $-523$     & $-526$     & $-0.6\%$ \\
Iron     & $(2) + (4{+}4) + (14) + (2)$                     & $-1260$    & $-1272$    & $-0.9\%$ \\
Krypton  & $(2) + (4{+}4) + (18) + (4{+}4)$                 & $-2766$    & $-2788$    & $-0.8\%$ \\
Xenon    & $(2) + (4{+}4) + (18) + (18) + (4{+}4)$          & $-7355$    & $-7438$    & $-1.1\%$ \\
Radon    & $(2) + (4{+}4) + (18) + (32) + (18) + (4{+}4)$   & $-22800$   & $-23560$   & $-3.2\%$ \\
\bottomrule
\end{tabular}
\end{table}

The pattern in Table~\ref{tab:atoms} is consistent across the range: light atoms (Carbon, Neon) come in slightly above the observed energy, suggesting that the Level-1 model is just slightly too tightly bound at low $Z$; mid-table atoms (Argon, Iron, Krypton, Xenon) come in slightly below, by about 1\%; the heaviest atom in the table (Radon) deviates by 3\% on the low side. We resist the reflex --- common in modern physics --- of attributing such a residual to ``relativistic effects.'' The electrons in RealQM are static charge densities, not moving particles, so a relativistic velocity correction has no referent here; and the total energy of a heavy atom is a large number built from many shells, computed by the simplified Atom Simulator with each sub-shell homogenised to spherical symmetry and the atom solved in reduced spatial detail. The few-percent deviation at $Z \sim 80$+ is the expected accumulation of \emph{these} simplifications --- the true RealQM total energy, computed in full spatial detail, is simply not what the reduced model returns --- and there is no warrant for charging it to relativity.

The full table for all elements Li--Rn is in the public Gallery (the Appendix). The pattern observed in Table~\ref{tab:atoms} extrapolates: across the working range Li--Xe the agreement is about 1\%, and for the heaviest atoms the deviation grows to a few percent as the spherically-symmetric simplification accumulates over more and more shells.

\section{What the agreement means}\label{sec:meaning}

We pause to spell out what the Level-1 result is and is not.

\paragraph{What it is.} Total atomic energies across Li--Rn are reproduced to about 1\% by a model with no fitted parameters. The shell structure $2, 8, 18, 32$ emerges from minimum-energy packing rather than being preset. The sub-shell radial split $(4{+}4)$ rather than $(8)$ is a prediction of the variational principle and is required for quantitative agreement. The Pauli exclusion principle does not appear as a separate postulate; the non-overlap constraint between electron territories does the same work.

This is a non-trivial validation of the central architectural choice of the framework: writing the $N$-electron wave function as a sum of one-electron wave functions on disjoint spatial domains. If that choice were wrong, the Level-1 model would not reproduce atomic energies at 1\% across the entire periodic table from a single principle. It does.

\paragraph{What it is not.} The Level-1 model does not capture every detail of atomic structure visible in StdQM. The angular distinction between $s$, $p$, $d$ orbitals within a shell is only partially represented; the $s/p$ splitting of an L-shell is replaced by a radial $(4{+}4)$ sub-shell split. Magnetic moments, fine structure, and hyperfine structure are not part of the model. The accuracy of 1\% is sufficient to anchor the higher levels of the hierarchy and to validate the multiphase formulation; it is not sufficient to compute, say, atomic ionisation potentials to spectroscopic precision.

The Level-1 model is in this sense the chemistry-relevant atomic backbone, not a complete atomic theory. It supplies what the rest of the hierarchy needs: an accurate ab initio anchor for the total atomic energy, against which the spherical homogenisation (Level 2) and the pseudo-kernel reduction (Level 3) are calibrated.

\section{Implications for the rest of Part III}\label{sec:to-part-iii}

The Level-1 agreement at 1\% across Li--Rn fixes the bottom of the hierarchy. From this point forward in the book, the higher-level architectures are derived by reduction from Level 1 and checked against observation in turn.

\begin{itemize}
\item Chapter~\ref{ch:h2} works out the H$_2$ covalent bond at parameter-free Level 1, with no kernel softening, and recovers the experimental dissociation energy to 3\%.
\item Chapter~\ref{ch:hydrides} sweeps Level-3 architectures across periodic-table groups, with kernel parameters $(Z_{\rm kernel}, r_c)$ calibrated against the Level-1 valence sector.
\item Chapter~\ref{ch:dimers} tests the resulting Level-3 model against the S66 H-bond benchmark and against reactive chemistry.
\item Chapter~\ref{ch:folding} stresses Level 4 with universal H-bond bias on polypeptide chains.
\item Chapter~\ref{ch:condensed} stresses Level 3 against three condensed-phase systems and quantifies what is missing.
\end{itemize}

The thread connecting these chapters is the Level-1 anchor of this chapter: each higher-level result rests on the Level-1 atomic energies agreeing with observation, and any discrepancy at a higher level must be traced either to the reduction or to physics not yet in the framework, not to the atomic anchor.

% --- End of Chapter 6 ---
